\backhead{FURTHER EXERCISES}{khred}
\begin{exlist}
\item Advance to the end of a border takes 90 min, the cut-off is at 240 min and recession is uniform in 15 min. With $Z=0.006\tau^{0.45}$ (m, min), find the depths at the head and the tail and the distribution uniformity (tail over mean of head and tail).\exans{72.6 mm; 59.7 mm; 0.90}
\item A furrow of 250 m receives 1.2 L\,s$^{-1}$ for 6 h; the furrow spacing is 0.75 m. Find the gross depth applied. If 55 mm is stored in the root zone, find the application efficiency.\exans{138 mm; 0.40}
\item Find the flow area, wetted perimeter and hydraulic radius of a furrow with bottom 0.15 m, side slope 1.5 and depth 0.1 m.\exans{0.030 m$^2$; 0.511 m; 0.059 m}
\item Using $Z=0.0045\tau^{0.5}$ (m, min), find the opportunity time needed for 75 mm, and the cut-off time if the advance takes 140 min and the required depth must be reached at the tail (neglect recession).\exans{278 min; 418 min}
\item A set lasts at most 48 hours. How many sets are needed to irrigate 60 furrows of 250 m with 1.5 L\,s$^{-1}$ each from a supply of 30 L\,s$^{-1}$, with 8 h per set?\exans{3 sets of 20 furrows; 24 h}
\end{exlist}
\backhead{ESSAY QUESTIONS}{khred}
{\small\itshape Answer each question in about one page, with equations, a sketch and at least one number from this chapter.}\par\vspace{3pt}
\begin{exlist}
\item Describe the phases of a surface irrigation (advance, wetting, depletion, recession) and how each affects uniformity.
\item Explain the volume-balance method of computing the advance, and its main assumption.
\item Discuss the choice between furrow, border and basin irrigation for a given soil and slope.
\item Explain why the intake of the soil is the most uncertain input of a surface design, and how SWAFL estimates it from the texture.
\item Discuss the trade-off between the application efficiency and the uniformity in surface irrigation.
\item Explain how cut-back and surge flow improve a furrow irrigation.
\end{exlist}
\dykband{EXERCISES IN SWAFL}
{\small Open SWAFL (desktop or online) with the benchmark project. The chapter of the \href{https://github.com/halilunaziru73-creator/Coupled-Water-Fluxes-Pore-Scale-Dynamics-Precision-Irrigation-Mulching-Feedback-Loops/raw/main/Software/v1.1.0/SWAFL_Desktop/docs/manual_v1.1.0/SWAFL_Manual_v1.1.0.pdf}{SWAFL manual} that explains each task is given in brackets; record your values in a table and compare them with what is expected (\emph{Check}).}\par\vspace{3pt}
\begin{exlist}
\item {\color{khblue}[Manual, Chapter 11]} Open \emph{Surface irrigation design} in \emph{Your soil} mode for the loam. Record the intake, the advance time, the cut-off and the efficiencies.\par\noindent{\small\itshape Check: the design table lists them.}
\item {\color{khblue}[Manual, Chapter 11]} Use \emph{Best inflow}; compare the efficiency at the best inflow with half and double that inflow.\par\noindent{\small\itshape Check: the efficiency is highest near the best inflow.}
\item {\color{khblue}[Manual, Chapter 11]} Read the tornado chart of the sensitivity. Which input changes the efficiency most?\par\noindent{\small\itshape Check: usually the intake or the inflow.}
\item {\color{khblue}[Manual, Chapter 11]} Read the irrigation schedule and the message \emph{When to irrigate}. Check the next date against the sensor and the trigger.\par\noindent{\small\itshape Check: the date follows the depletion.}
\end{exlist}
